% ============================================================================
%  ESTADÍSTICA DESCRIPTIVA — Sesión 4
%  Semana 4: 31 de agosto al 5 de septiembre de 2026
%  1.7 Métodos gráficos: histogramas y polígonos de frecuencia
%  Duración prevista: 3 horas
% ============================================================================
\documentclass[aspectratio=169,10pt]{beamer}

\usepackage[utf8]{inputenc}
\IfFileExists{spanish.ldf}{\usepackage[spanish,es-noshorthands]{babel}}{}
\usepackage{amsmath,amssymb}
\usepackage{booktabs}
\usepackage{tikz}
\usepackage{listings}
\usepackage{pgfplots}
\pgfplotsset{compat=1.16}
\usetikzlibrary{patterns}

\usetheme{Uniminuto}

\uniasignatura{Estadística descriptiva}
\uniprofesor{Andrés Felipe Puerta Velez}
\unicorreo{andres.puerta-v@uniminuto.edu.co}
\uninrc{CBASBL031 - CBASBL141}
\unisemana{Sesión 4 — Semana 4 · 31 de agosto al 5 de septiembre de 2026}

\newcommand{\listatiny}{\setbeamerfont{itemize/enumerate body}{size=\tiny}\tiny}

% Estilo común de los gráficos
\pgfplotsset{
  uniplot/.style={
    width=8.4cm, height=3.9cm,
    axis lines=left,
    tick label style={font=\tiny,color=uniGris},
    label style={font=\tiny,color=uniGris},
    title style={font=\scriptsize\bfseries,color=uniAzul},
    every axis plot/.append style={fill=uniAzul!70,draw=uniAzul},
    ymajorgrids, grid style={uniGris!25,very thin},
  }
}

\newcommand{\fx}[1]{{\ttfamily\color{uniAzul}#1}}

% ------------------------------------------------------------- CÓDIGO R ----
\lstdefinestyle{Rcode}{
  language=R,
  basicstyle=\ttfamily\scriptsize,
  keywordstyle=\color{uniAzul}\bfseries,
  commentstyle=\color{uniGris}\itshape,
  stringstyle=\color{uniTitulo},
  showstringspaces=false, numbers=none, breaklines=true,
  columns=fullflexible,
  backgroundcolor=\color{uniCremaClaro},
  frame=leftline, framerule=1.2pt, rulecolor=\color{uniNaranja},
  xleftmargin=0.35cm, framexleftmargin=0.35cm,
  framexrightmargin=0.15cm, framextopmargin=2pt, framexbottommargin=2pt,
  aboveskip=0.15cm, belowskip=0.15cm
}
\lstdefinestyle{Rmini}{
  language=R,
  basicstyle=\ttfamily\tiny,
  keywordstyle=\color{uniAzul}\bfseries,
  commentstyle=\color{uniGris}\itshape,
  stringstyle=\color{uniTitulo},
  showstringspaces=false, numbers=none, breaklines=true,
  columns=fullflexible,
  backgroundcolor=\color{uniCremaClaro},
  frame=leftline, framerule=1.2pt, rulecolor=\color{uniNaranja},
  xleftmargin=0.35cm, framexleftmargin=0.35cm,
  framexrightmargin=0.15cm, framextopmargin=2pt, framexbottommargin=2pt,
  aboveskip=0.12cm, belowskip=0.12cm
}
\lstdefinestyle{Rout}{
  basicstyle=\ttfamily\tiny\color{black},
  showstringspaces=false, numbers=none, breaklines=true,
  columns=fullflexible,
  backgroundcolor=\color{uniGris!10},
  frame=leftline, framerule=1.2pt, rulecolor=\color{uniGris},
  xleftmargin=0.35cm, framexleftmargin=0.35cm,
  framexrightmargin=0.15cm, framextopmargin=2pt, framexbottommargin=2pt,
  aboveskip=0.12cm, belowskip=0.12cm
}

\begin{document}

\portada

\begin{frame}{Orden del día}
  \begin{enumerate}
    \item Llamado a lista y revisión de los ejercicios propuestos
    \item Repaso: de los datos crudos a la tabla
    \item \azul{Gráficos para variables cualitativas: barras y circular}
    \item \azul{Gráficos para variables cuantitativas: histograma, polígono y ojiva}
    \item Taller guiado con el caso Almacén Bello
    \item Formas de la distribución y errores frecuentes
  \end{enumerate}

  \vspace{0.25cm}
  {\footnotesize\color{uniGris} Sesión de 3 horas. Traiga la tabla de la sesión
  3: hoy la convertimos en imagen.}
\end{frame}

\begin{frame}{Repaso: dónde quedamos}
  En la sesión 3 construimos tablas de frecuencia para el caso Almacén Bello.

  \vspace{0.2cm}
  \begin{alertblock}{La pregunta de hoy}
    \footnotesize
    Una tabla con seis filas y siete columnas es correcta, pero
    \destacar{nadie la mira en una junta}. ¿Cómo convertimos esa tabla en una
    imagen que muestre de un vistazo dónde se concentran los datos?
  \end{alertblock}

  \vspace{0.2cm}
  \begin{block}{La regla que gobierna toda la sesión}
    \footnotesize
    \azul{El tipo de variable decide el gráfico.} Barras y circular para
    cualitativas; histograma, polígono y ojiva para cuantitativas agrupadas.
    Usar un histograma para una variable nominal es un error, no una
    preferencia estética.
  \end{block}
\end{frame}

\begin{frame}{El mapa de los gráficos}
  \scriptsize
  \begin{center}
  \begin{tabular}{@{}p{3.0cm}p{3.4cm}p{5.4cm}@{}}
    \toprule
    \textbf{Gráfico} & \textbf{Sirve para} & \textbf{Qué muestra} \\
    \midrule
    Barras          & Cualitativa (nominal u ordinal) &
      Compara categorías; barras \azul{separadas} \\
    Circular        & Cualitativa con pocas categorías &
      Composición del total (partes de 100\%) \\
    Histograma      & Cuantitativa agrupada &
      Forma de la distribución; barras \azul{pegadas} \\
    Polígono        & Cuantitativa agrupada &
      Perfil de la distribución; permite comparar grupos \\
    Ojiva           & Cuantitativa agrupada &
      Frecuencia acumulada; lectura de percentiles \\
    \bottomrule
  \end{tabular}
  \end{center}

  \vspace{0.1cm}
  \begin{alertblock}{Barras separadas frente a barras pegadas}
    \scriptsize
    En un \destacar{diagrama de barras} las categorías son etiquetas
    independientes: las barras van separadas. En un \destacar{histograma} el
    eje es una recta numérica continua y las barras se tocan.
  \end{alertblock}
\end{frame}

\seccion{Gráficos para variables cualitativas}

\begin{frame}{Diagrama de barras — canal de compra}
  \centering
  \begin{tikzpicture}
    \begin{axis}[uniplot,
      title={Canal de compra de 40 clientes — Almacén Bello},
      ybar, bar width=0.85cm,
      ylabel={Frecuencia absoluta $f_i$},
      symbolic x coords={Tienda,App,Web,Teléfono},
      xtick=data, ymin=0, ymax=16, enlarge x limits=0.18,
      nodes near coords, nodes near coords style={font=\tiny,color=uniAzul}]
      \addplot coordinates {(Tienda,14) (App,13) (Web,9) (Teléfono,4)};
    \end{axis}
  \end{tikzpicture}

  \vspace{-0.1cm}
  \begin{alertblock}{Cuatro reglas de construcción}
    \scriptsize
    (1) El eje vertical \destacar{empieza en cero}, siempre. \quad
    (2) Todas las barras tienen el \destacar{mismo ancho}. \quad
    (3) Las barras van \destacar{separadas}. \quad
    (4) Si la variable es nominal, se ordena de mayor a menor frecuencia; si es
    ordinal, se respeta el orden natural.
  \end{alertblock}
\end{frame}

\begin{frame}{Diagrama circular — el cálculo de los ángulos}
  \footnotesize
  \begin{columns}[T,onlytextwidth]
    \begin{column}{0.50\textwidth}
      \centering\scriptsize
      \begin{tabular}{@{}lrrr@{}}
        \toprule
        \textbf{Canal} & \textbf{$f_i$} & \textbf{$\%_i$} &
        \textbf{Ángulo} \\
        \midrule
        Tienda   & 14 & 35{,}0 & $126^\circ$ \\
        App      & 13 & 32{,}5 & $117^\circ$ \\
        Web      &  9 & 22{,}5 & $81^\circ$  \\
        Teléfono &  4 & 10{,}0 & $36^\circ$  \\
        \midrule
        \textbf{Total} & \textbf{40} & \textbf{100} & $\mathbf{360^\circ}$ \\
        \bottomrule
      \end{tabular}

      \vspace{0.15cm}
      \[ \text{ángulo}_i = h_i \times 360^\circ \]
      \[ \tfrac{14}{40}\times 360 = 126^\circ \]
    \end{column}
    \begin{column}{0.46\textwidth}
      \centering
      \begin{tikzpicture}[scale=0.95]
        \def\r{1.75}
        \fill[uniAzul]     (0,0) -- (90:\r) arc (90:-36:\r) -- cycle;
        \fill[uniNaranja]  (0,0) -- (-36:\r) arc (-36:-153:\r) -- cycle;
        \fill[uniAmarillo] (0,0) -- (-153:\r) arc (-153:-234:\r) -- cycle;
        \fill[uniGris]     (0,0) -- (-234:\r) arc (-234:-270:\r) -- cycle;
        \draw[white,thick] (0,0)--(90:\r); \draw[white,thick] (0,0)--(-36:\r);
        \draw[white,thick] (0,0)--(-153:\r); \draw[white,thick] (0,0)--(-234:\r);
        \node[font=\tiny,white] at (27:1.05) {35\%};
        \node[font=\tiny,white] at (-94:1.05) {32,5\%};
        \node[font=\tiny,white] at (166:1.05) {22,5\%};
        \node[font=\tiny,white] at (108:1.15) {10\%};
        \node[font=\tiny,right,color=uniGris] at (2.0,0.75) {Tienda};
        \node[font=\tiny,right,color=uniGris] at (2.0,0.35) {App};
        \node[font=\tiny,right,color=uniGris] at (2.0,-0.05) {Web};
        \node[font=\tiny,right,color=uniGris] at (2.0,-0.45) {Teléfono};
        \fill[uniAzul]    (1.75,0.65) rectangle (1.95,0.85);
        \fill[uniNaranja] (1.75,0.25) rectangle (1.95,0.45);
        \fill[uniAmarillo](1.75,-0.15) rectangle (1.95,0.05);
        \fill[uniGris]    (1.75,-0.55) rectangle (1.95,-0.35);
      \end{tikzpicture}
    \end{column}
  \end{columns}

  \vspace{0.1cm}
  \begin{alertblock}{Cuándo NO usar el circular}
    \scriptsize
    Con \destacar{más de 5 o 6 categorías} los sectores se vuelven
    indistinguibles, y el ojo humano compara mal ángulos. Con muchas
    categorías, o cuando lo importante es comparar, \azul{use barras}.
  \end{alertblock}
\end{frame}

\begin{frame}{Ejercicio 1 — en clase}
  \begin{alertblock}{En parejas (15 minutos)}
    \footnotesize
    Retome la tabla de satisfacción de la sesión 3 ($n = 30$):

    \vspace{0.15cm}
    \centering\scriptsize
    \begin{tabular}{@{}lrrrrr@{}}
      \toprule
      \textbf{Satisfacción} & Muy bajo & Bajo & Medio & Alto & Muy alto \\
      \midrule
      \textbf{$f_i$} & 3 & 5 & 10 & 8 & 4 \\
      \bottomrule
    \end{tabular}

    \vspace{0.25cm}
    \raggedright
    \begin{enumerate}\setlength{\itemsep}{-0.05cm}
      \item Calcule $h_i$, $\%_i$ y el ángulo de cada sector.
      \item Dibuje el diagrama de barras. ¿En qué orden van las categorías?
      \item ¿Es apropiado un diagrama circular aquí? Justifique.
      \item Escriba dos frases de interpretación para el gerente.
    \end{enumerate}
  \end{alertblock}
\end{frame}

\begin{frame}{Ejercicio 1 — solución}
  \footnotesize
  \begin{columns}[T,onlytextwidth]
    \begin{column}{0.44\textwidth}
      \centering\scriptsize
      \begin{tabular}{@{}lrrr@{}}
        \toprule
        \textbf{Nivel} & \textbf{$f_i$} & \textbf{$\%_i$} & \textbf{Áng.} \\
        \midrule
        Muy bajo & 3 & 10{,}0 & $36^\circ$ \\
        Bajo     & 5 & 16{,}7 & $60^\circ$ \\
        Medio    & 10 & 33{,}3 & $120^\circ$ \\
        Alto     & 8 & 26{,}7 & $96^\circ$ \\
        Muy alto & 4 & 13{,}3 & $48^\circ$ \\
        \midrule
        \textbf{Tot.} & \textbf{30} & \textbf{100} & $\mathbf{360^\circ}$ \\
        \bottomrule
      \end{tabular}
    \end{column}
    \begin{column}{0.53\textwidth}
      \begin{exampleblock}{Puntos 2 y 3}
        \scriptsize
        \textbf{Orden:} de muy bajo a muy alto. La variable es \azul{ordinal},
        así que ordenar por frecuencia \destacar{destruiría la información}
        que el gráfico debe mostrar.

        \vspace{0.1cm}
        \textbf{Circular:} es aceptable (5 categorías), pero las barras son
        mejores porque conservan visualmente el orden de la escala.
      \end{exampleblock}
    \end{column}
  \end{columns}

  \vspace{0.1cm}
  \begin{alertblock}{Punto 4: dos frases posibles}
    \scriptsize
    ``El \destacar{60,0\%} de los clientes reporta satisfacción media o
    inferior.'' \quad ``Solo 1 de cada 6 clientes está en el nivel más alto:
    hay margen de mejora en el servicio.''
  \end{alertblock}
\end{frame}

\seccion{Gráficos para variables cuantitativas}

\begin{frame}{Histograma — caso Almacén Bello}
  \centering
  \begin{tikzpicture}
    \begin{axis}[uniplot,
      title={Gasto mensual de 40 clientes (miles de pesos)},
      ybar interval, bar width=1,
      xlabel={Gasto (miles \$)}, ylabel={$f_i$},
      xtick={112,152,192,232,272,312,352},
      ymin=0, ymax=12.5, x tick label style={font=\tiny},
      xticklabel style={/pgf/number format/1000 sep=}]
      \addplot coordinates {(112,4) (152,7) (192,11) (232,9) (272,6) (312,3) (352,0)};
    \end{axis}
  \end{tikzpicture}

  \vspace{-0.1cm}
  \begin{alertblock}{Cómo se construye}
    \scriptsize
    En el eje horizontal van los \destacar{límites de clase} (no las
    categorías); en el vertical, la frecuencia. Cada barra ocupa exactamente el
    ancho de su clase, y las barras \destacar{se tocan}: el gasto es una
    variable continua y no hay huecos entre 152 y 152.
  \end{alertblock}
\end{frame}

\begin{frame}{Polígono de frecuencias}
  \centering
  \begin{tikzpicture}
    \begin{axis}[uniplot,
      title={Polígono de frecuencias — gasto mensual},
      xlabel={Marca de clase $x_i$}, ylabel={$f_i$},
      xtick={92,132,172,212,252,292,332,372},
      ymin=0, ymax=12.5, xmin=80, xmax=384,
      every axis plot/.append style={fill=none}]
      \addplot[color=uniTitulo,thick,mark=*,mark size=1.6pt,
               mark options={fill=uniTitulo}]
        coordinates {(92,0) (132,4) (172,7) (212,11) (252,9) (292,6) (332,3) (372,0)};
    \end{axis}
  \end{tikzpicture}

  \vspace{-0.1cm}
  \begin{alertblock}{Los dos puntos que casi todos olvidan}
    \scriptsize
    El polígono une las \destacar{marcas de clase} a la altura de $f_i$, y
    debe \destacar{cerrarse sobre el eje} con una clase ficticia \azul{antes}
    de la primera ($x=92$, $f=0$) y otra \azul{después} de la última
    ($x=372$, $f=0$).
  \end{alertblock}
\end{frame}

\begin{frame}{Ojiva — frecuencias acumuladas}
  \centering
  \begin{tikzpicture}
    \begin{axis}[uniplot,
      title={Ojiva — gasto mensual acumulado},
      xlabel={Gasto (miles \$)}, ylabel={$F_i$},
      xtick={112,152,192,232,272,312,352},
      ymin=0, ymax=44, xmin=108, xmax=356,
      every axis plot/.append style={fill=none}]
      \addplot[color=uniAzul,thick,mark=*,mark size=1.6pt,
               mark options={fill=uniAzul}]
        coordinates {(112,0) (152,4) (192,11) (232,22) (272,31) (312,37) (352,40)};
      \draw[uniTitulo,dashed] (axis cs:108,20) -- (axis cs:225,20)
                              -- (axis cs:225,0);
    \end{axis}
  \end{tikzpicture}

  \vspace{-0.1cm}
  \begin{alertblock}{Para qué sirve}
    \scriptsize
    La ojiva se traza sobre los \destacar{límites superiores} de clase, nunca
    sobre las marcas. ¿Cuál es el gasto por debajo del cual está la mitad de
    los clientes? Se busca $F=20$, se corta la curva y se baja:
    \azul{alrededor de 225 mil}. Eso es la mediana.
  \end{alertblock}
\end{frame}

\begin{frame}{Los tres gráficos vienen de la misma tabla}
  \scriptsize
  \begin{center}
  \begin{tabular}{@{}lrrrl@{}}
    \toprule
    \textbf{Clase} & \textbf{$x_i$} & \textbf{$f_i$} & \textbf{$F_i$} &
    \textbf{Se usa en} \\
    \midrule
    $[112,\ 152)$ & 132 &  4 &  4 & Histograma, polígono y ojiva \\
    $[152,\ 192)$ & 172 &  7 & 11 & \\
    $[192,\ 232)$ & 212 & 11 & 22 & Clase modal: la barra más alta \\
    $[232,\ 272)$ & 252 &  9 & 31 & \\
    $[272,\ 312)$ & 292 &  6 & 37 & \\
    $[312,\ 352]$ & 332 &  3 & 40 & La ojiva termina en $n=40$ \\
    \bottomrule
  \end{tabular}
  \end{center}

  \vspace{0.05cm}
  \begin{center}\scriptsize
  \begin{tabular}{@{}lll@{}}
    \toprule
    \textbf{Gráfico} & \textbf{Eje horizontal} & \textbf{Eje vertical} \\
    \midrule
    Histograma & Límites de clase & $f_i$ (o $h_i$) \\
    Polígono   & Marcas de clase $x_i$ & $f_i$ (o $h_i$) \\
    Ojiva      & Límites superiores & $F_i$ (o $H_i$) \\
    \bottomrule
  \end{tabular}
  \end{center}
\end{frame}

\begin{frame}{Ejercicio 2 — en clase (taller principal)}
  \begin{alertblock}{En grupos de 3 (25 minutos)}
    \scriptsize
    Retome la tabla del tiempo de espera de la sesión 3 ($n=30$):

    \vspace{0.1cm}
    \centering\scriptsize
    \begin{tabular}{@{}lrrrrrr@{}}
      \toprule
      \textbf{Clase} & $[4,9)$ & $[9,14)$ & $[14,19)$ & $[19,24)$ & $[24,29)$ &
      $[29,34]$ \\
      \midrule
      $x_i$ & 6{,}5 & 11{,}5 & 16{,}5 & 21{,}5 & 26{,}5 & 31{,}5 \\
      $f_i$ & 4 & 6 & 6 & 5 & 5 & 4 \\
      \bottomrule
    \end{tabular}

    \vspace{0.15cm}
    \raggedright
    \begin{enumerate}\setlength{\itemsep}{-0.1cm}
      \item Dibuje el histograma en papel milimetrado o cuadriculado.
      \item Superponga el polígono de frecuencias, con las clases ficticias.
      \item Complete $F_i$ y dibuje la ojiva.
      \item Lea en la ojiva el tiempo por debajo del cual está el 50\% de los
            clientes.
      \item Compare la forma de este histograma con el del gasto: ¿cuál está
            más concentrado?
    \end{enumerate}
  \end{alertblock}
\end{frame}

\begin{frame}{Ejercicio 2 — solución (histograma y polígono)}
  \centering
  \begin{tikzpicture}
    \begin{axis}[uniplot, width=7.2cm, height=3.6cm,
      title={Tiempo de espera en caja (minutos)},
      ybar interval, bar width=1,
      xlabel={Minutos}, ylabel={$f_i$},
      xtick={4,9,14,19,24,29,34}, ymin=0, ymax=7.5]
      \addplot coordinates {(4,4) (9,6) (14,6) (19,5) (24,5) (29,4) (34,0)};
      \addplot[color=uniTitulo,thick,mark=*,mark size=1.3pt,fill=none,
               mark options={fill=uniTitulo}]
        coordinates {(1.5,0) (6.5,4) (11.5,6) (16.5,6) (21.5,5) (26.5,5)
                     (31.5,4) (36.5,0)};
    \end{axis}
  \end{tikzpicture}

  \vspace{-0.1cm}
  \begin{alertblock}{Punto 5: comparación de formas}
    \scriptsize
    El histograma del gasto tiene un \destacar{pico claro} en $[192,232)$ y
    cae a ambos lados: los datos están \azul{concentrados}. El del tiempo de
    espera es casi \destacar{plano}: todos los tiempos son igual de probables,
    los datos están \azul{dispersos}. En la unidad 2 pondremos un número a esa
    diferencia (la desviación estándar).
  \end{alertblock}
\end{frame}

\begin{frame}{Ejercicio 2 — solución (ojiva)}
  \footnotesize
  \begin{columns}[T,onlytextwidth]
    \begin{column}{0.36\textwidth}
      \centering\scriptsize
      \begin{tabular}{@{}lrr@{}}
        \toprule
        \textbf{Clase} & \textbf{$f_i$} & \textbf{$F_i$} \\
        \midrule
        $[4,9)$   & 4 &  4 \\
        $[9,14)$  & 6 & 10 \\
        $[14,19)$ & 6 & 16 \\
        $[19,24)$ & 5 & 21 \\
        $[24,29)$ & 5 & 26 \\
        $[29,34]$ & 4 & 30 \\
        \bottomrule
      \end{tabular}
    \end{column}
    \begin{column}{0.62\textwidth}
      \centering
      \begin{tikzpicture}
        \begin{axis}[uniplot, width=6.4cm, height=3.9cm,
          xlabel={Minutos}, ylabel={$F_i$},
          xtick={4,9,14,19,24,29,34}, ymin=0, ymax=33, xmin=3, xmax=35,
          every axis plot/.append style={fill=none}]
          \addplot[color=uniAzul,thick,mark=*,mark size=1.3pt,
                   mark options={fill=uniAzul}]
            coordinates {(4,0) (9,4) (14,10) (19,16) (24,21) (29,26) (34,30)};
          \draw[uniTitulo,dashed] (axis cs:3,15) -- (axis cs:18.2,15)
                                  -- (axis cs:18.2,0);
        \end{axis}
      \end{tikzpicture}
    \end{column}
  \end{columns}

  \vspace{0.05cm}
  \begin{alertblock}{Punto 4}
    \scriptsize
    Se busca $F = n/2 = 15$ en el eje vertical y se corta la ojiva:
    el resultado está \destacar{alrededor de 18 minutos}. Con los datos
    originales la mediana es $(17+18)/2 = 17{,}5$: la lectura gráfica es
    aproximada, pero sirve.
  \end{alertblock}
\end{frame}

\seccion{Formas de la distribución}

\begin{frame}{Cuatro formas que hay que reconocer a simple vista}
  \centering
  \begin{tikzpicture}[scale=0.86]
    \begin{scope}
      \foreach \x/\h in {0/0.4,0.5/0.9,1/1.5,1.5/1.9,2/1.5,2.5/0.9,3/0.4}
        \fill[uniAzul!70] (\x,0) rectangle (\x+0.5,\h);
      \draw[uniGris] (0,0)--(3.5,0);
      \node[font=\tiny,color=uniAzul] at (1.75,-0.35) {Simétrica};
    \end{scope}
    \begin{scope}[xshift=4.2cm]
      \foreach \x/\h in {0/1.9,0.5/1.6,1/1.1,1.5/0.8,2/0.5,2.5/0.3,3/0.15}
        \fill[uniNaranja!80] (\x,0) rectangle (\x+0.5,\h);
      \draw[uniGris] (0,0)--(3.5,0);
      \node[font=\tiny,color=uniAzul,align=center] at (1.75,-0.45)
        {Asimétrica positiva\\(cola a la derecha)};
    \end{scope}
    \begin{scope}[xshift=8.4cm]
      \foreach \x/\h in {0/0.15,0.5/0.3,1/0.5,1.5/0.8,2/1.1,2.5/1.6,3/1.9}
        \fill[uniNaranja!80] (\x,0) rectangle (\x+0.5,\h);
      \draw[uniGris] (0,0)--(3.5,0);
      \node[font=\tiny,color=uniAzul,align=center] at (1.75,-0.45)
        {Asimétrica negativa\\(cola a la izquierda)};
    \end{scope}
    \begin{scope}[xshift=12.6cm]
      \foreach \x/\h in {0/1.6,0.5/1.0,1/0.5,1.5/0.4,2/0.6,2.5/1.2,3/1.7}
        \fill[uniGris!60] (\x,0) rectangle (\x+0.5,\h);
      \draw[uniGris] (0,0)--(3.5,0);
      \node[font=\tiny,color=uniAzul] at (1.75,-0.35) {Bimodal};
    \end{scope}
  \end{tikzpicture}

  \vspace{0.25cm}
  \begin{alertblock}{Por qué importa}
    \scriptsize
    La forma decide \destacar{qué medida resume mejor}. En una distribución
    simétrica la media es un buen resumen. En una asimétrica positiva —típica
    de salarios y gastos— la media se va detrás de la cola y la
    \azul{mediana} describe mejor al caso típico. Una distribución
    \azul{bimodal} suele significar que hay \destacar{dos grupos mezclados}.
  \end{alertblock}
\end{frame}

\begin{frame}{¿Qué forma tiene el gasto del Almacén Bello?}
  \footnotesize
  \begin{columns}[T,onlytextwidth]
    \begin{column}{0.52\textwidth}
      \centering
      \begin{tikzpicture}
        \begin{axis}[uniplot, width=6.6cm, height=3.7cm,
          ybar interval, bar width=1,
          xtick={112,192,272,352}, ymin=0, ymax=12.5,
          xlabel={Gasto (miles \$)}, ylabel={$f_i$}]
          \addplot coordinates {(112,4) (152,7) (192,11) (232,9) (272,6)
                                (312,3) (352,0)};
        \end{axis}
      \end{tikzpicture}
    \end{column}
    \begin{column}{0.45\textwidth}
      \begin{exampleblock}{Lectura}
        \scriptsize
        Pico en $[192,232)$, caída rápida a la izquierda y \azul{cola larga a
        la derecha}: es una distribución \destacar{asimétrica positiva}.

        \vspace{0.15cm}
        Tiene sentido económico: hay un piso de gasto (nadie gasta menos de
        cero) pero no hay techo.
      \end{exampleblock}
    \end{column}
  \end{columns}

  \vspace{0.1cm}
  \begin{alertblock}{La consecuencia práctica}
    \scriptsize
    Si el gerente reporta el \destacar{gasto promedio}, ese número estará
    inflado por los pocos clientes de la cola. Para describir al cliente
    típico conviene la \azul{mediana}. Esta es la discusión central de la
    unidad 2.
  \end{alertblock}
\end{frame}

\begin{frame}{Errores frecuentes en gráficos}
  \scriptsize
  \begin{center}
  \begin{tabular}{@{}p{5.4cm}p{6.4cm}@{}}
    \toprule
    \textbf{Error} & \textbf{Por qué es un error} \\
    \midrule
    Eje vertical que no empieza en cero &
      Exagera diferencias pequeñas: es la manipulación gráfica más común \\
    Histograma con barras separadas &
      Sugiere huecos en una variable continua \\
    Barras pegadas en variable nominal &
      Sugiere continuidad donde solo hay etiquetas \\
    Gráfico circular con 12 categorías &
      Los sectores dejan de ser comparables \\
    Efectos 3D y sombras &
      Distorsionan las áreas; no aportan información \\
    Polígono sin cerrar sobre el eje &
      El área deja de coincidir con la del histograma \\
    Ojiva trazada sobre marcas de clase &
      Se traza sobre \azul{límites superiores}: $F_i$ acumula hasta ahí \\
    Gráfico sin título, sin ejes rotulados o sin fuente &
      El lector no puede saber qué está viendo \\
    \bottomrule
  \end{tabular}
  \end{center}
\end{frame}

\begin{frame}[fragile]{Los mismos gráficos, en R}
  \scriptsize
\begin{lstlisting}[style=Rmini]
# Cualitativa: barras separadas y circular
barplot(table(canal), col = "steelblue",
        main = "Canal de compra de 40 clientes", ylab = "Frecuencia")
pie(table(canal), main = "Participacion por canal")

# Histograma con NUESTROS limites de clase
limites <- seq(112, 352, by = 40)
hist(gasto, breaks = limites, right = FALSE, col = "steelblue",
     main = "Gasto mensual de 40 clientes", xlab = "Miles de pesos")

# Poligono: marcas de clase, cerrado sobre el eje
xi <- head(limites, -1) + 20
lines(c(92, xi, 372), c(0, as.vector(table(clases)), 0), type = "b", col = "tomato")

# Ojiva: acumulada sobre los limites superiores
plot(limites, c(0, cumsum(table(clases))), type = "b",
     main = "Ojiva del gasto", xlab = "Miles de pesos", ylab = "Fi")
\end{lstlisting}

  \vspace{-0.1cm}
  \begin{alertblock}{El argumento que no se puede olvidar}
    \tiny
    Sin \fx{breaks = limites}, \fx{hist()} elige sus propias clases con la regla
    de Sturges y \destacar{el gráfico deja de coincidir con la tabla}. En R,
    como a mano, las clases las decide usted.
  \end{alertblock}
\end{frame}

\begin{frame}{Cierre de la sesión}
  \begin{block}{Lo que debe quedar claro hoy}
    \scriptsize
    \begin{itemize}\setlength{\itemsep}{0.02cm}
      \item Cualitativas: \azul{barras} (separadas) y \azul{circular}
            (ángulo $= h_i \times 360^\circ$).
      \item Cuantitativas agrupadas: \azul{histograma} (barras pegadas sobre
            los límites), \azul{polígono} (marcas de clase, cerrado sobre el
            eje) y \azul{ojiva} (límites superiores, $F_i$).
      \item Los tres salen de la \destacar{misma tabla}: cambia qué columna se
            grafica contra qué.
      \item La \azul{forma} de la distribución anticipa qué medida resumirá
            mejor los datos.
    \end{itemize}
  \end{block}

  \vspace{0.1cm}
  \begin{alertblock}{Trabajo independiente}
    \scriptsize
    Resolver los 20 ejercicios de la lámina siguiente. En la sesión 5
    integramos todo el proceso —de los datos crudos al informe— y hacemos el
    repaso para el examen 1.
  \end{alertblock}
\end{frame}

\begin{frame}{Ejercicios propuestos}
  \vspace{-0.1cm}
  \begin{beamercolorbox}[wd=\textwidth,sep=0.10cm]{block body}
    \centering\scriptsize
    \textbf{1--8:} elección del gráfico. \quad
    \textbf{9--14:} cálculo. \quad
    \textbf{15--20:} construcción e interpretación.
  \end{beamercolorbox}
  \listatiny
  \vspace{0.15cm}
  \begin{columns}[T,onlytextwidth]
    \begin{column}{0.47\textwidth}
      \begin{enumerate}\setlength{\itemsep}{0.16cm}
        \item Gráfico adecuado para: ciudad de residencia
        \item Para: gasto mensual agrupado en 6 clases
        \item Para: nivel educativo alcanzado
        \item Para: la proporción acumulada de clientes
        \item ¿Barras separadas o pegadas en un histograma?
        \item ¿Sobre qué valores se traza el polígono?
        \item ¿Sobre qué valores se traza la ojiva?
        \item ¿Por qué el eje $y$ debe empezar en cero?
        \item $h_i = 0{,}22$: ¿cuál es el ángulo del sector?
        \item Ángulo $= 54^\circ$: ¿cuál es $\%_i$?
      \end{enumerate}
    \end{column}
    \begin{column}{0.50\textwidth}
      \begin{enumerate}\setlength{\itemsep}{0.16cm}
        \setcounter{enumi}{10}
        \item Clases $[10,20)\dots[50,60)$: ¿cuáles son las clases ficticias
              del polígono?
        \item $n=40$ y la ojiva termina en\dots
        \item Si $f_i=11$ y $n=40$: altura de la barra en escala $h_i$
        \item ¿Cambia la forma del histograma si grafico $h_i$ en vez de $f_i$?
        \item Dibuje el histograma de: 3, 8, 12, 9, 5, 3
        \item ¿Cuál es la clase modal en ese histograma?
        \item ¿Qué forma tiene esa distribución?
        \item Con $F_i$: 3, 11, 23, 32, 37, 40. Dibuje la ojiva
        \item Lea en esa ojiva el valor con $F=20$
        \item Un gráfico muestra ventas de 98 a 102 con eje de 97 a 103.
              ¿Qué problema tiene?
      \end{enumerate}
    \end{column}
  \end{columns}
\end{frame}

\begin{frame}{Respuestas}{Autoevaluación de los ejercicios propuestos}
  \listatiny
  \vspace{0.15cm}
  \begin{columns}[T,onlytextwidth]
    \begin{column}{0.47\textwidth}
      \begin{enumerate}\setlength{\itemsep}{0.16cm}
        \item Barras (nominal); circular si son pocas ciudades
        \item Histograma y polígono
        \item Barras en orden natural (ordinal)
        \item Ojiva
        \item Pegadas
        \item Sobre las marcas de clase $x_i$
        \item Sobre los límites superiores de clase
        \item Porque truncarlo exagera las diferencias
        \item $0{,}22\times360 = 79{,}2^\circ$
        \item $54/360 = 0{,}15 \Rightarrow 15\%$
      \end{enumerate}
    \end{column}
    \begin{column}{0.50\textwidth}
      \begin{enumerate}\setlength{\itemsep}{0.16cm}
        \setcounter{enumi}{10}
        \item Marcas ficticias en 5 y 65, con $f=0$
        \item En 40, es decir en $n$
        \item $h_i = 11/40 = 0{,}275$
        \item No: solo cambia la escala del eje vertical
        \item Barras de altura 3, 8, 12, 9, 5, 3, pegadas
        \item La tercera clase ($f=12$)
        \item Asimétrica positiva (cola a la derecha)
        \item Puntos crecientes que terminan en 40
        \item Entre la tercera y la cuarta clase
        \item Eje truncado: exagera diferencias mínimas
      \end{enumerate}
    \end{column}
  \end{columns}
\end{frame}

\begin{frame}{Referencias}
  \footnotesize
  \begin{itemize}
    \item[{[1]}] Monroy Saldívar, S. (2008). \textit{Estadística descriptiva}.
      Instituto Politécnico Nacional.
      \texttt{https://elibro.net/es/lc/uniminuto/titulos/74722}
    \item[{[2]}] Ordóñez Fernández, F. F. \& González Fernández, J. (2021).
      \textit{Estadística descriptiva paso a paso}. Ediciones Pirámide.
      \texttt{https://elibro.net/es/lc/uniminuto/titulos/215449}
    \item[{[3]}] Llinás Solano, H. (2014). \textit{Introducción a la
      estadística con aplicaciones en ciencias sociales}. Universidad del
      Norte. \texttt{https://elibro.net/es/lc/uniminuto/titulos/70062}
  \end{itemize}

  \vspace{0.35cm}
  \begin{center}
    {\Large\color{uniAzul}\textbf{¿Preguntas?}}
  \end{center}
\end{frame}

\end{document}
